% !TEX root = main.tex
\section{서론}
\subsection{제품정보의 신뢰성과 공급망 근거}
디지털 제품 여권(Digital Product Passport, DPP)은 제품의 특성, 환경정보, 이력을 해당 제품과 연결하여 필요한 주체가 조회할 수 있도록 하는 체계다.
EU 배터리 규정은 2027년 2월 18일부터 시장에 출시되거나 사용이 개시되는 전기차용 배터리, 경량 운송수단용 배터리, 용량 2 kWh 초과 산업용 배터리에 전자 배터리 여권을 요구한다.
여권의 환경정보에는 탄소발자국과 재활용 원료 함량이 포함되며, 각 정보 의무는 해당 조항에서 정한 일정에 따라 적용된다. [1]

최종 제조사는 이러한 정보를 제공하기 위해 공급망의 앞선 단계에서 생성된 자료와 증빙을 확보하고 관리해야 한다.
배터리 규정 제39조는 셀·모듈 공급자가 제조자의 규정 준수에 필요한 정보와 문서를 제공하도록 한다.
제77조 제4항은 책임 사업자가 여권 정보의 정확성, 완전성, 최신성을 확보하도록 요구한다. [1]

그러나 증빙을 확보하는 것만으로 후속 제품정보의 타당성이 보장되지는 않는다.
원료를 여러 제품으로 나누거나 다른 원료와 함께 가공하면, 그 원료에 관한 회계 수치를 각 제품에 어떻게 배분하고 반영했는지도 확인해야 한다.
\textbf{증빙이 유효하게 발급되었는지와 그 증빙을 사용한 후속 계산이 올바른지는 서로 다른 검증 대상이다.}
본 연구는 생산·이전 과정에서 사용된 회계 자료가 최종 제품의 지속가능성 주장에 반영되기까지의 관계를 검증하고 관리하는 문제를 다룬다. [2]

이러한 검증은 민감한 공급망 정보를 보호하면서 이루어져야 한다.
배터리 규정은 공개 정보와 접근이 제한된 정보를 구분하고, 데이터의 무결성·보안·프라이버시를 요구한다.
Babel 등도 공급망 정보가 통제되지 않은 채 확산될 수 있다는 우려가 DPP 정보 공유의 장애가 된다고 지적한다. [1, 2]
따라서 본 연구는 공개할 제품정보와 이를 뒷받침하는 상세 거래자료를 구분한다.
기업은 회계 수치와 어떤 원료를 공급받아 어떤 제품을 생산했는지에 관한 연결을 일반 검증자에게 직접 공개하지 않는다.
대신 해당 자료를 사용한 거래와 최종 주장 발급이 정해진 규칙을 따랐음을 검증받는다.
\textbf{거래의 유효성을 검증하는 일과 상세 공급망 이력을 열람하는 일을 서로 다른 권한으로 구분하는 것이다.}

\subsection{정상 검증을 보완하는 사후 대응}
원천 증빙이 타당한지는 외부의 공급·검사·감사 절차에서 판단한다.
반면 원장의 거래 검증은 수락된 자료를 사용한 계산·배분·이전이 정해진 규칙을 따랐는지 확인한다.
공급자의 증빙을 바탕으로 생산과 이전이 진행된 뒤, 추가 검사에서 그 증빙의 오류가 발견될 수 있다.
\textbf{자료에 오류가 있더라도 그 자료를 사용한 계산은 규칙에 맞을 수 있다.}
따라서 정상 거래를 검증하는 것과 별도로, 기존 자료를 계속 근거로 인정할지 재평가하고 그 판단을 이후 사용에 반영할 절차가 필요하다.

권한 있는 감사자가 특정 근거의 추가 사용을 중단하기로 결정했다면, 그 근거에서 파생되어 현재 남아 있는 기록에도 결정이 적용되어야 한다.
이미 여러 기업에서 생산과 이전이 진행됐다면, 최초 기록에만 표시를 남겨서는 후속 거래나 새로운 주장 발급에 쓰이는 기록을 제어할 수 없다.
비공개로 유지하던 거래 연결을 조사해 관련 기록을 찾고, 그 기록을 소비하는 다음 거래가 거부되도록 해야 한다.
근거의 재평가와 사용 중단 결정은 감사자가 맡는다.
시스템은 원장에 실제로 수락된 이력을 조사하고, 승인된 사용 제한을 집행하도록 지원한다.
이러한 사후 대응은 정상 거래 검증을 보완하는 안전장치다.

본 연구의 질문은 다음과 같다.
\textbf{민감한 공급망 정보를 일반 검증자에게 직접 공개하지 않으면서, 최종 제품의 주장을 그 회계 근거와 연결하고, 해당 근거에 대한 사후 판단을 후속 기록의 사용에 반영하려면 어떤 기록·검증·권한 구조가 필요한가?}

\subsection{선행연구와 본 연구의 과제}
기존 연구는 비공개 검증과 공급망 정보 관리에 필요한 기반을 제공한다.
Babel 등은 생산 단계의 자격증명을 연결하고 원료 지분의 이전과 소비를 관리하는 DPP 구성을 제시했다.
Udokwu와 Craß는 자격증명으로 확인한 입력을 사용해 지속가능성 주장을 증명하고 검증하는 체계를 다룬다.
Sahai 등은 분할·합류·가공을 포함하는 공급망 문서의 관계와 추적을 다루며, 입력 문서의 참조를 원장에 기록한다.
비공개 자산 거래 프로토콜인 CAP(Configurable Asset Privacy)는 정책에 따른 열람·동결과 실제 거래에 결합된 감사 자료를 지원한다. [2, 3, 4, 5]

본 연구는 이러한 기법을 사용해 \textbf{DPP Claim의 발급 근거 확인과 사후 사용 통제를 같은 공급망 이력에 연결하는 시스템}을 구체화한다.
최종 제품을 검증할 때에는 Claim이 유효한 회계 상태와 발급 조건에 근거했는지 확인해야 한다.
문제가 있는 출처에서 조사를 시작할 때에는 그 근거가 어떤 후속 기록으로 이어졌고, 현재 어떤 기록의 사용을 제한해야 하는지 확인해야 한다.

두 업무가 같은 이력을 다루려면 네 대상이 서로 일치해야 한다.
거래에서 검증한 근거, Claim 발급에 사용한 상태, 감사자가 복원한 입력·출력 관계, 원장이 소비를 허용하거나 거부하는 대상이다.
\textbf{본 연구의 과제는 이 대응 관계와 이를 유지하는 권한·처리 규칙을 명확히 정의하는 것이다.}
기존 비공개 거래 기법을 공급망에 적용하더라도 이러한 관계와 규칙은 별도로 정해야 한다.

\subsection{zkDPP의 설계와 기여}
zkDPP는 생산·분할·가공·기업 간 이전을 공유 원장에 기록되는 상태 변화로 표현한다.
참여자는 거래의 실제 입력과 출력이 회계·소유권·권한 규칙을 만족함을 영지식증명으로 보인다.
원장은 이 증명과 입력 기록의 현재 사용 가능 상태를 함께 검사한다.

최종 Claim 발급에서는 회계 상태가 발급 정책을 충족하는지 검증한 뒤 해당 기록을 소비한다.
그 결과로 지정된 제품정보와 정책에 연결된 공개 참조인 \textbf{DPP Claim}을 등록한다.
외부 검증자는 상세 거래자료를 열람하지 않고도 DPP Claim이 해당 제품정보와 원장의 발급 이력에 연결되어 있는지 확인한다.

각 거래는 사후 추적에 필요한 실제 입력 연결과 이후 소비 식별 정보를 암호화해 남긴다.
거래 증명은 이 정보가 실제 입력·출력과 일치하는지도 검증한다.
감사자는 이를 복호화하여 거래 당시 검증된 연결을 따라 이력을 조사한다.

감사 자료를 복호화하는 데이터 키는 별도의 임계 암호화로 보호한다.
$N$명의 위원 중 $K$명 이상의 협력과 외부 승인 절차를 거쳐 이 데이터 키를 복구하고 감사자에게 전달한다.
$K\ge2$이고 초기 키 사본이 올바르게 관리된다는 조건에서, 위원 한 명의 비밀 조각만으로는 감사 자료를 열 수 없다.
승인된 감사자는 전달받은 키로 기록의 연결을 로컬에서 조사하고, 후속 조치의 판단과 집행을 담당한다.

사용 중단이 승인되면 감사자는 확정된 원장 시점을 기준으로 문제 출처의 후속 기록을 추적한다.
그중 아직 소비되지 않은 기록 전체를 동결 대상으로 삼으며, 완료되지 않은 기업 간 이전도 포함한다.
동결 상태인 기록은 후속 가공·이전이나 새로운 Claim 발급에 사용할 수 없다.

발급과 동결의 결과는 공유 원장이 두 요청을 수락한 순서와 동일한 사용 상태에 따라 결정된다.
동결이 먼저 확정되면 차단이 유지되는 동안 해당 기록의 Claim 발급을 거부한다.
발급이 먼저 확정되면 그 결과를 이미 완료된 발급 이력으로 남긴다.
이렇게 하면 여러 갈래로 나뉜 공급망에서도 \textbf{이미 발급된 Claim을 확인하는 일과 아직 소비되지 않은 기록의 다음 사용을 막는 일을 구분}할 수 있다.

전체 차단의 완료 여부는 처음 정한 대상 목록을 유지한 채 하나의 공통 확정 시점에서 판단한다.
대상이 존재하고, 그 전체가 같은 시점에 여전히 미소비·차단 상태이며, 처리 결과에도 불명확한 항목이 없어야 완료로 보고한다.
동결 전에 대상이 소비되었거나 조회·제출 결과가 불명확하면 해당 대상과 사유를 결과에 남긴다.
동결할 대상이 없는 경우도 별도로 보고한다.
따라서 일부 동결 성공을 전체 대상의 차단 완료로 해석하지 않는다.

본 논문의 기여는 다음 두 가지다.
\begin{enumerate}
    \item \textbf{비공개 회계 검증, DPP Claim 발급, 승인된 사후 대응을 연결하는 시스템 설계.}
          생산과 이전에 따라 회계 근거가 어떻게 변하는지 검증하고, 이를 최종 Claim에 연결한다.
          감사키 접근과 후속 기록의 사용 제한까지 포함하는 프로토콜을 정의하며, 같은 기록에 대한 발급과 동결, 이미 발급된 Claim과 남은 기록의 처리를 일관된 상태 변화·권한 규칙으로 명세한다.
          이를 통해 구현에 필요한 증명, 원장 검사, 감사 절차의 관계를 제시한다.
    \item \textbf{시스템의 보장과 실패 조건에 대한 조건부 분석.}
          명시한 암호·원장·권한 가정 아래 정확성, 건전성, 프라이버시를 정의하고 증명한다.
          거래·Claim·감사 자료가 실제 생성된 기록과 일치함을 보이고, 이를 정확한 추적과 고정된 전체 대상의 소비 차단으로 연결한다.
          전체 대상이 공통 확인 시점에 미소비·차단 상태이면 차단이 유지되는 동안 추가 소비가 거부됨을 보인다.
          자료 오류나 소비와 동결의 경합으로 완료할 수 없는 경우도 구분한다.
          프라이버시 분석은 참여자 비밀과 감사·상태 정보가 공개되지 않은 정상 이력을 수동 외부 관찰자가 보는 경우에 한정한다.
\end{enumerate}
